
\documentclass[landscape]{article}
\usepackage[a4paper,margin=1cm]{geometry}
\usepackage{multicol}
\usepackage{pdflscape}
\usepackage{amsmath,amssymb,amsthm}
\usepackage{tcolorbox}
\tcbuselibrary{skins,breakable}
\setlength{\columnsep}{0.8cm}

\begin{document}
\begin{landscape}
\begin{multicols}{4}

% -------------------- CLASE 3 --------------------
Clase 3

Ya estamos en condiciones de definir uno de los conceptos más importantes asociados a los polinomios: el de cero o raíz de un polinomio. El cálculo de los ceros de un polinomio será de suma utilidad para resolver ecuaciones, inecuaciones y cualquier tipo de problema que requiera la factorización de polinomios.
\begin{tcolorbox}[colback=blue!5!white,colframe=blue!75!black,title=\textbf{Definición (Cero o raíz de un polinomio)}]
Decimos que un número real $\alpha$ es un cero o raíz del polinomio $p(x)$ si $p(\alpha) = 0$.
\end{tcolorbox}
\noindent\textbf{Ejemplo 1.10.} Calcule el valor numérico del polinomio $p(x) = -x^3 + 4x - 5$ para $x = 3$.
\noindent\textbf{Solución:}
$p(3) = -(3)^3 + 4(3) - 5 = -27 + 12 - 5 = -20$.
\noindent\textbf{Ejemplo 1.11.} Verifique que $-2$ es una raíz del polinomio $p(x) = x^3 + x^2 - 5x - 6$.
\noindent\textbf{Solución:}
$p(-2) = (-2)^3 + (-2)^2 - 5(-2) - 6 = -8 + 4 + 10 - 6 = 0$,
por lo que $p(-2)=0$ confirma que $x=-2$ es una raíz.
\noindent\textbf{Ejemplo 1.12.} Determine el valor de $\alpha$ para que $4$ sea raíz del polinomio $p(x) = -x^3 + \alpha x^2 + 2x + 6$.
\noindent\textbf{Solución:} Para que $4$ sea una raíz de $p(x)$, debe ocurrir $p(4) = 0$. Luego:
\begin{align*}
0 &= -(4)^3 + \alpha (4)^2 + 2(4) + 6,\\
0 &= -64 + 16\alpha + 8 + 6,\\
0 &= -50 + 16\alpha,\\
16\alpha &= 50,\\
\alpha &= \frac{50}{16} = \frac{25}{8}.
\end{align*}
Recordemos que un polinomio $p(x)$ es divisible por otro polinomio $q(x)$ si el resto de dividir $p$ por $q$ es igual a cero. El siguiente teorema muestra que cuando el divisor es de la forma $q(x) = x - k$, podemos determinar si $p(x)$ es divisible por $q(x)$ sin necesidad de realizar la división: basta con calcular $p(k)$.
\begin{tcolorbox}[colback=red!5!white,colframe=red!75!black,title=\textbf{Teorema del resto}]
El resto de dividir $p(x)$ por un polinomio de la forma $x - k$ (con $k$ un número real cualquiera) es igual a $p(k)$.

\medskip\noindent\textbf{Demostración:}
Si $p(x)$ es un polinomio de grado mayor o igual que 1, por el Algoritmo de la División sabemos que existen polinomios $c(x)$ y $r(x)$ tales que:
\[p(x) = (x - k)c(x) + r(x),\]
donde $r(x)$ es el resto y cumple $\deg r < \deg(x - k) = 1$. En particular, $r(x)$ es un polinomio constante (o $r(x)=0$ si la división es exacta). Evaluando en $x = k$ se tiene:
\[p(k) = (k - k)c(k) + r(k) = r,\]
de donde $p(k) = r$. Es decir, el resto $r$ coincide con $p(k)$, lo que prueba el teorema.
\end{tcolorbox}

\noindent\textbf{Ejemplo 1.13.} Calcule el resto de dividir $p(x) = -x^4 - 2x^3 + 3x^2 + x + 1$ por $q(x) = x + 1$.

\noindent\textbf{Solución:}
$p(-1) = -(-1)^4 - 2(-1)^3 + 3(-1)^2 + (-1) + 1 = -1 + 2 + 3 - 1 + 1 = 4$.
Por lo tanto, el resto de la división es $r(x) = 4$.

\noindent\textbf{Ejemplo 1.14.} Halle el valor de $a \in \mathbb{R}$ para que, al dividir $p(x) = x^4 + 3x^3 + ax^2 - 13$ por $x + 3$, el resto sea igual a $5$.

\noindent\textbf{Solución:}
Usando el teorema del resto, se tiene $p(-3) = 5$. Entonces:
\[
(-3)^4 + 3(-3)^3 + a(-3)^2 - 13 = 5,
\]
de donde $9a = 18$ y por lo tanto $a = 2$.

Del Teorema del Resto se deduce el siguiente resultado:

\begin{tcolorbox}[colback=red!5!white,colframe=red!75!black,title=\textbf{Teorema}]
Si $\alpha \in \mathbb{R}$ es una raíz del polinomio $p(x)$, entonces $p(x)$ es divisible por $x - \alpha$.

\medskip\noindent\textbf{Demostración:}
Por el Teorema del Resto, el resto de dividir $p(x)$ por $x - \alpha$ es $p(\alpha)$. Pero $p(\alpha) = 0$, luego el resto es 0 y se concluye que $x - \alpha$ divide a $p(x)$.
\end{tcolorbox}

% -------------------- more content would follow --------------------

\columnbreak

% Placeholder for Clase 4 y 5 (incluir resto del contenido aquí si se desea)

\end{multicols}
\end{landscape}
\end{document}
